\documentclass[11pt,a4paper]{article}

\usepackage{hyperref}
\usepackage{geometry}
\usepackage{titlesec}
\usepackage{setspace}

\geometry{margin=1in}
\setstretch{1.2}

\titleformat{\section}{\large\bfseries}{\thesection}{1em}{}

\title{Using the LaTeX Build Container}
\author{Henrique Moura}
% \date{}

\begin{document}

\maketitle

\section{Overview}

This repository provides a fully containerized LaTeX build environment based on Ubuntu and a complete TeX Live installation. The container automatically compiles the document located in the \texttt{doc/} directory when started. Bibliography files are supported, and optional LaTeX packages can be installed dynamically using an \texttt{install.lst} file.

\section{Cloning the Repository}

To obtain the repository locally, run:

\begin{verbatim}
git clone git@github.com:h3dema/docker-latex.git
cd docker-latex
\end{verbatim}


\section{Building the Container}

The project includes a \texttt{docker-compose.yaml} file that builds the image and sets up the workspace mount.

\begin{verbatim}
docker-compose build
\end{verbatim}

This step installs:
\begin{itemize}
    \item Ubuntu base system
    \item Full TeX Live distribution
    \item \texttt{latexmk} for automated compilation
    \item A custom entrypoint script that handles package installation and PDF generation
\end{itemize}

\section{Running the Container}

To compile the LaTeX project inside \texttt{doc/}, run:

\begin{verbatim}
docker-compose up
\end{verbatim}

The container will:

\begin{enumerate}
    \item Check for \texttt{/workspace/doc/install.lst}.  
          If present, each package listed (one per line) is installed using \texttt{tlmgr}.
    \item Determine the main \texttt{.tex} file.  
          The default is \texttt{main.tex}, but it can be overridden:
\begin{verbatim}
LATEX_MAIN=thesis.tex docker-compose up
\end{verbatim}
    \item Compile the document using \texttt{latexmk}, including bibliography support.
\end{enumerate}

The resulting PDF will appear inside your local \texttt{doc/} directory.

\begin{itemize}
    \item The \texttt{install.lst} file is optional and allows dynamic installation of additional LaTeX packages.
    \item Bibliography files (\texttt{.bib}) are automatically handled by \texttt{latexmk}.
    \item The container always compiles the document on startup; no manual commands inside the container are required.
\end{itemize}

\end{document}
